% R038：三条候选线采用同一坐标范围、同一单位长度；
% 数据为 (1,2),(2,3),(3,5),(4,6)

\begin{center}
	\begin{tikzpicture}[
			x=1cm,
			y=1cm,
			>=Stealth,
			every node/.style={font=\small},
			axis/.style={->,black!70,thick},
			data/.style={
					circle,
					fill=blue!80!black,
					inner sep=2.3pt
				},
			res/.style={red!78!black,thick}
		]

		\path[use as bounding box] (0,0) rectangle (12.4,7.65);

		\node[
			font=\large\bfseries,
			text=CExplanation
		] at (6.2,7.38)
		{同一坐标尺度下比较：残差平方和最小的候选线最优};

		% 三幅图统一使用相同的横纵坐标缩放
		\def\xstep{0.66}
		\def\ystep{0.50}
		\def\axisY{1.82}

		\foreach \o/\b/\a/\q/\eq/\title in {
		0/1.0/1.5/1.00/{$\hat y=x+1.5$}/{候选},
		4.2/1.4/0.5/0.20/{$\hat y=1.4x+0.5$}/{最优},
		8.4/1.8/-0.5/1.00/{$\hat y=1.8x-0.5$}/{候选}
		}{

		% 卡片背景
		\draw[
			rounded corners=5pt,
			fill=blue!2,
			draw=CExplanation,
			line width=.8pt
		]
		(\o+.08,.48) rectangle (\o+3.92,6.88);

		% 突出中间的最优候选线
		\ifdim\o pt=4.2pt
			\draw[
				orange!85!black,
				line width=1.5pt,
				rounded corners=5pt
			]
			(\o+.08,.48) rectangle (\o+3.92,6.88);
		\fi

		% 卡片标题
		\node[
			font=\bfseries,
			text=CExplanation
		] at (\o+2,6.57)
		{\title：$b=\b$};

		% 回归直线方程
		\node[
			font=\normalsize,
			text=CExplanation
		] at (\o+2,6.08)
		{\eq};

		% x 轴
		\draw[axis]
		(\o+.62,\axisY)
		--
		(\o+3.48,\axisY)
		node[above right=-1pt] {$x$};

		% y 轴
		\draw[axis]
		(\o+.62,\axisY)
		--
		(\o+.62,5.55)
		node[above] {$y$};

		% x 轴刻度
		\foreach \xx in {1,2,3,4}{
				\pgfmathsetmacro{\px}{\o+.62+\xstep*\xx}

				\draw
				(\px,\axisY-.08)
				--
				(\px,\axisY+.08);

				\node[below=3pt] at (\px,\axisY)
				{$\xx$};
			}

		% y 轴刻度
		\foreach \yy in {1,2,3,4,5,6,7}{
				\pgfmathsetmacro{\py}{\axisY+\ystep*\yy}

				\draw
				(\o+.54,\py)
				--
				(\o+.70,\py);

				\node[left=2pt] at (\o+.62,\py)
				{$\yy$};
			}

		% 四个原始数据点
		\foreach \xx/\yy in {1/2,2/3,3/5,4/6}{
				\pgfmathsetmacro{\px}{\o+.62+\xstep*\xx}
				\pgfmathsetmacro{\py}{\axisY+\ystep*\yy}

				\node[data] at (\px,\py) {};
			}

		% 候选回归线的两个端点
		\pgfmathsetmacro{\xone}{\o+.62+\xstep}
		\pgfmathsetmacro{\yone}{\axisY+\ystep*(\a+\b)}

		\pgfmathsetmacro{\xfour}{\o+.62+4*\xstep}
		\pgfmathsetmacro{\yfour}{\axisY+\ystep*(\a+4*\b)}

		% 绘制候选回归线
		\draw[
			CStatement,
			very thick
		]
		(\xone,\yone)
		--
		(\xfour,\yfour);

		% 绘制纵向残差
		\foreach \xx/\yy in {1/2,2/3,3/5,4/6}{
				\pgfmathsetmacro{\px}{\o+.62+\xstep*\xx}
				\pgfmathsetmacro{\py}{\axisY+\ystep*\yy}
				\pgfmathsetmacro{\ly}{\axisY+\ystep*(\a+\b*\xx)}

				\draw[res]
				(\px,\py)
				--
				(\px,\ly);
			}

		% 样本中心点
		\pgfmathsetmacro{\cx}{\o+.62+2.5*\xstep}
		\pgfmathsetmacro{\cy}{\axisY+4*\ystep}

		\fill[orange!90!black]
		(\cx,\cy) circle (3.1pt);

		\node[
			anchor=south,
			xshift=-10pt,
			text=orange!85!black
		] at (\cx,\cy)
		{$(2.5,4)$};

		% 残差平方和：单独留出底部空间
		\node[
			font=\bfseries\small,
			text=red!78!black
		] at (\o+2,.91)
		{$Q=\sum e_i^2=\q$};
		}

		% 图下注释
		\node[
			text=red!78!black,
			anchor=west,
			font=\small
		] at (.22,.14)
		{每条红线段均为纵向残差；只画三条候选线，严格全局最优见后文。};

	\end{tikzpicture}
\end{center}
